\subsection{Os seis zeros de uma função}
\label{sec:divisors-seis-zeros}
\index{multiplicidade@multiplicidade!interpolacao local@interpolação local}

Que função escreverá Alice para a representação de três pelas cinco
parcelas de vinte e dois? Conhecemos a sua forma geral e sabemos que
existe. Podemos agora descobrir-lhe os coeficientes. A conta tem uma
particularidade que a recta do exemplo binário ainda deixava na sombra.
Das seis ocorrências do divisor, três pertencem ao mesmo ponto. Para
construir a função é preciso traduzir cada uma delas numa condição
sobre os coeficientes. Comecemos por dispor os dados pela ordem em
que os vamos usar.
\begin{center}
\begin{tabular}{lccc}
 Lugar & Coordenadas & Ocorrências & Proveniência \\ \hline
 \(4P\) & \((3,1)\) & 3 & tabela \\
 \(2P\) & \((6,3)\) & 1 & tabela \\
 \(8P\) & \((13,7)\) & 1 & tabela \\
 \(-3P\) & \((10,11)\) & 1 & inverso de \(W\)
\end{tabular}
\end{center}
A última distinção interessa. O ponto \(-3P\) coincide nesta curva
com \(16P\), o quinto lugar da tabela. O vector de multiplicidades
atribuiu zero a esse lugar. A ocorrência que ali vemos veio da
parcela \([-W]\) de \eqref{eq:divisors-fixed-size-divisor}.
Havemos de reencontrá-la na verificação, com coeficiente um.

Tomemos \(t=x-3\). Nas vizinhanças de \((3,1)\), a equação da
curva determina \(y\) como série em \(t\), pois \(2y=2\) admite
inversa em \(\mathbb F_{17}\). Podemos descobrir os termos de que
precisamos por mera comparação. A substituição de \(3+t\) em
\(x^3+2x+2\) dá
\[
                 1+12t+9t^2+t^3\pmod {17}.
\]
Escrevamos \(y=1+ut+vt^2+wt^3+\cdots\). No quadrado desta
expressão, os coeficientes dos três primeiros graus positivos são
\(2u\), \(u^2+2v\) e \(2uv+2w\). A primeira igualdade dá
\(u=6\). Na segunda, \(36+2v=9\) dá \(v=12\).
Na terceira, \(144+2w=1\) dá \(w=5\). Fica, assim,
\begin{equation}
              y=1+6t+12t^2+5t^3+\cdots .
 \label{eq:divisors-six-local}
\end{equation}
A ordenada e o declive apareceram nos dois primeiros termos. O
terceiro, que descreve a diferença entre a curva e a sua tangente,
é indispensável sempre que se quer impor a terceira ocorrência.

A translação \(x=3+t\) conserva os limites de grau de \(f_0\)
e \(f_1\). Por isso, a função pode escrever-se como uma combinação
de \(1,t,t^2,t^3,y,ty\). Para obter um zero pelo menos triplo,
basta agora conservar, de cada combinação, as que começam no grau
três. Tomemos as três combinações
\[
 t^3,\qquad
 U_0=y-1-6t-12t^2,\qquad
 U_1=t(y-1-6t).
\]
Em \(U_0\) subtraímos os três primeiros termos da série de \(y\).
Em \(U_1\) o próprio \(t\) já deu um grau, pelo que basta
subtrair dois. Assim se encontra toda a família procurada. Com
efeito, fixados os coeficientes de \(t^3,y,ty\), os coeficientes
de \(1,t,t^2\) ficam determinados, sucessivamente, pela anulação
dos três primeiros termos. A família escreve-se
\[
                       f=k t^3+c U_0+e U_1.
\]
Nos seis lugares contados com multiplicidade tem de se esgotar um
pólo de ordem seis. Entre os monómios permitidos, só \(t^3\) tem
esse peso. Logo, para a função que procuramos, \(k\ne0\). Nesta
fase é cómodo fixar a escala pondo \(k=1\).

Restam três lugares simples. Em cada um podemos avaliar directamente
\(t^3,U_0,U_1\), usando a ordenada que está no quadro. Por exemplo,
em \((6,3)\) temos \(t=3\), \(t^3=27=10\),
\(U_0=3-1-18-108=12\) e \(U_1=3(3-1-18)=3\),
tudo módulo dezassete. As três linhas dão
\[
\begin{array}{c|rrr|l}
 &t^3&U_0&U_1& f=0\\ \hline
 (6,3)   &10&12&3 & 10+12c+3e=0\\
 (13,7)  &14& 4&4 & 14+4c+4e=0\\
 (10,11) & 3& 9&14& 3+9c+14e=0
\end{array}
\]
Da segunda linha sai \(c+e=5\), pois \(-14=3\) e
\(4^{-1}=13\). Substituir \(e=5-c\) na primeira deixa
\(25+9c=0\), isto é, \(9c=9\). Encontram-se \(c=1,e=4\).
A última linha vale \(3+9+56=68=0\). Esta última concordância
tinha uma razão geométrica antes de ganhar números. Fixados cinco
zeros de uma função com pólo de ordem seis, a lei de grupo
determina o lugar do sexto. As três equações do quadro exprimem,
na linguagem dos coeficientes, aquela mesma dependência.

Voltemos à abscissa \(x\). A expressão \(t^3+U_0+4U_1\) fica
\begin{equation}
 f(x,y)=1+12x+6x^2+x^3+(6+4x)y.
 \label{eq:divisors-six-polynomial}
\end{equation}
Pode seguir-se, por exemplo, a origem do coeficiente de \(y\).
Ele era \(1+4t\) e, ao substituir \(t=x-3\), passa a
\(1-12+4x=6+4x\). A parte sem \(y\) era
\(t^3-36t^2-10t-1\). A expansão de \((x-3)^3\) e de
\((x-3)^2\), com a mesma redução, dá os quatro primeiros termos
da equação. A soma de pontos transformou-se enfim numa lista
de seis números que Alice pode guardar.

Há duas maneiras instrutivas de reler estes coeficientes. A primeira
consiste em regressar à expansão local. Em \(4P\), o primeiro
termo de \(U_0\) é \(5t^3\), e o de \(U_1\) é \(12t^3\).
O coeficiente de \(t^3\) em \(f\) vale
\(1+5+4\cdot12=54=3\). O zero tem exactamente ordem três.
Nos outros três lugares, o quadro impôs pelo menos uma ocorrência.
Chegámos a seis, e a parcela \(x^3\) fixa o pólo de ordem seis.
Ficam explicados todos os zeros da função sobre o fecho algébrico,
incluindo a sua multiplicidade.

A segunda leitura parte da norma do capítulo. Conservemos a
convenção \(N_f=f_0^2-(x^3+2x+2)f_1^2\). Ao multiplicar
\((f_0+y f_1)(f_0-y f_1)\), os termos ímpares em \(y\)
cancelam-se e a substituição da equação da curva dá exactamente
esse polinómio. Com os coeficientes de \eqref{eq:divisors-six-polynomial},
vem
\begin{align*}
 N_f(x)
 &=x^6+13x^5+12x^4+10x^3+11x^2+9x+14\\
 &=(x-3)^3(x-6)(x-10)(x-13)
              \qquad\text{em }\mathbb F_{17}[x].
\end{align*}
A última escrita pode obter-se logo pela geometria. Cada zero de
\(f\) fornece à norma a sua abscissa, com a multiplicidade
correspondente quando o outro ponto dessa fibra é regular e tem
valor não nulo. Conhecemos seis ocorrências e o coeficiente de
\(x^6\) é um. A multiplicação acima permite também percorrer a
conta no sentido inverso, começando pelo polinómio que Alice tem
na mão.

Recuperar as abscissas da norma ainda deixa, para cada uma, a
pergunta pela ordenada. Avaliemos \(f\) nos dois pontos da curva
que têm a mesma abscissa.
\[
\begin{array}{c|rr|rr}
 x& y& f(x,y)&-y&f(x,-y)\\ \hline
 3&1&0&16&15\\
 6&3&0&14&7\\
 10&11&0&6&8\\
 13&7&0&10&4
\end{array}
\]
Por exemplo, para \(x=10\), a parcela que acompanha \(y\) é
\(6+40=12\), a parcela constante vale \(4\), e
\(4+12\cdot11=136=0\). Com a ordenada oposta obtém-se
\(4+12\cdot6=76=8\). A norma guarda a abscissa de \(-3P\)
e \(f\) guarda ainda o sinal que a soma de pontos exigiu.
Também na fibra de \(4P\), a ordem três pertence inteiramente ao
ponto de ordenada um. Na figura os círculos vazios indicam o outro
ponto de cada fibra. A ordem das ordenadas é a dos representantes
\(0,\ldots,16\) do corpo.

\begin{center}
\begin{tikzpicture}[x=.48cm,y=.195cm,font=\small]
 \draw[->] (1.8,0)--(15.4,0) node[right] {\(x\)};
 \draw[->] (1.8,0)--(1.8,18.5) node[above] {\(y\)};
 \foreach \x/\ya/\yb in {3/1/16,6/3/14,10/11/6,13/7/10} {
   \draw[gray!60,densely dotted] (\x,0)--(\x,16.7);
   \draw (\x,0)--(\x,-.5) node[below] {\(\x\)};
   \fill (\x,\ya) circle[radius=2pt];
   \draw[fill=white] (\x,\yb) circle[radius=2pt];
 }
 \draw[thin] (3,1.55)--(3.38,4.4);
 \node[anchor=south west,inner sep=1pt] at (3.38,4.4) {\(4P\ (\times3)\)};
 \node[anchor=west] at (6.12,3) {\(2P\)};
 \node[anchor=east] at (9.85,11) {\(-3P\)};
 \node[anchor=west] at (13.1,7) {\(8P\)};
\end{tikzpicture}
\end{center}

Para a prova que fixa o coeficiente de \(x\), falta escolher a
escala correspondente. Esse coeficiente em \(f\) é doze, cuja
inversa vale dez. Multiplicar cada termo por dez deixa
\begin{equation}
 \widehat f=10+x+9x^2+10x^3+(9+6x)y .
 \label{eq:divisors-six-normalized}
\end{equation}
Na construção contámos seis ocorrências. Escrevemos as coordenadas
no corpo de dezassete elementos. Na relação \(W=3P\),
o coeficiente de \(P\) está tomado módulo dezanove. A lista
normalizada contém cinco coeficientes livres e um coeficiente fixo.
As cinco multiplicidades são outra lista, também com cinco
posições, mas com um encargo diferente. Os coeficientes livres
permitem avaliar uma função. As multiplicidades descrevem quais os
pontos que devem responder por essas avaliações. Na
secção~\ref{sec:gbp-sete-portas} seguiremos ambas as listas até às
mesmas portas de um circuito.
